% The record and the ring. Left: one event record, byte offsets outside the
% column. Right: where the working buffers live in SRAM.
\begin{tikzpicture}[font=\sffamily\scriptsize]

% ---------------- left column: one event record, low offsets at the bottom
\node[font=\sffamily\bfseries\small, anchor=west] at (0mm,52mm) {One event record};
\node[font=\sffamily\tiny, anchor=west, text width=34mm] at (0mm,47mm)
  {500 ms window at 1667 Hz, six bytes per sample};
\memregion{0}{9}{mmext}{record header\\{\tiny 16 B, carries both scale codes}}{0}
\memregion{9}{12}{mmram}{pre-trigger samples\\{\tiny 167 x 6 B = 1002 B}}{16}
\memregion{21}{14}{mmram}{post-trigger samples\\{\tiny 667 x 6 B = 4002 B}}{1018}
\memregion{35}{8}{mmperiph}{pressure trace\\{\tiny 13 x 5 B = 65 B}}{5020}
\node[mmaddr] at (-0.5mm,43mm) {5085};

% ---------------- right column: SRAM placement
\begin{scope}[xshift=58mm]
\node[font=\sffamily\bfseries\small, anchor=west] at (0mm,52mm) {Where it lives in SRAM};
\node[font=\sffamily\tiny, anchor=west, text width=34mm] at (0mm,47mm)
  {AXI SRAM at 0x2400\,0000, 1024 kB in three banks. Budget only};
\memregion{0}{9}{mmres}{drivers stack and console}{}
\memregion{9}{14}{mmram}{pre-trigger ring\\{\tiny 2048 x 6 B = 12288 B}}{}
\memregion{23}{12}{mmram}{record staging\\{\tiny one open window}}{}
\memregion{35}{8}{mmres}{free}{}
\node[mmaddr, anchor=west] at (35mm,4.5mm) {shared};
\node[mmaddr, anchor=west] at (35mm,16mm) {12 kB};
\node[mmaddr, anchor=west] at (35mm,29mm) {8 kB};
\node[mmaddr, anchor=west] at (35mm,39mm) {the rest};
\end{scope}

\node[umlnote, text width=96mm, anchor=north west] at (0mm,-6mm)
  {The ring exists so that the part of the event the fine channel could not
   detect in time is already stored when the coarse channel says the event
   happened. Which SRAM these buffers are placed in is a separate decision with
   its own consequences for the transfer engines; that is Chapter 19.};
\end{tikzpicture}
